%================================================================
% CHAPTER 5: SYSTEM TESTING
%================================================================
\chapter{SYSTEM TESTING}

\section{Introduction}

This chapter includes the results which are discussed briefly by running the system during testing. It is error free and all the features work properly. The main objective is to find the errors and try to sort out bugs and errors for increasing the efficiency of the system. After completing the testing process the next target is to analyze the results. Outputs are examined after the process. The next process is the validation process, which contains results that are analyzed and then resolves bugs, and then implements on the system.

\section{Testing and Validation}

The actual meaning of testing is to search out the errors which occur in the system and to make sure that the system should be error free at every stage. Sometimes there are many errors which cannot easily be found without the testing process. Testing is used to make the system better because in testing many of the errors become solved and this increases the reliability and accuracy of the system. The main purpose of testing is to satisfy the user, and this is the first priority. So, to make the system validation reliable, the testing checks all the components of the system one by one and removes all the errors during the process. When we test the first component and the result is not correct, it causes some errors in the first phase; we then test that component again and again until we achieve the desired output. The same process must be done with all other components to make sure the availability of right information. Another major aspect of the system is to check the validation of the user, whether the user enters some wrong input into the system. For example, if a user does not select or fill in the search field then an error is generated due to the incompletion of the process, which is where validation occurs.

\subsection{System Testing}

The aim of the system remains the same; it does not change no matter which method is used to test all the limits of the proposed system during the process. Our main purpose is to make sure that the system is reliable and works correctly and properly and gives correct results as the user wants.

\subsection{Unit Testing}

In this process of testing, the system is divided into components and then tested one by one. It is an engineering approach by which errors can be easily found. It also helps to make the performance of the system very good. The components of the system are checked one by one because it is not certain which component the error is in.

\subsection{Integration Testing}

Integration testing checks the errors and the working of the system, and whether components work properly or not, after the completion of unit testing.

\subsection{User Testing}

In this testing, the user decides about the system how much the system is reliable and easy to understand. The user checks the whole system very deeply because the user is the main person to use the system. After checking the whole system, the user gives feedback about the system. If users are not satisfied with the work and give negative remarks about the system, then those parts are integrated again that the user might not like. Making the system user-friendly is of primary importance because if the user rejects the system performance then all the effort and work is wasted.

\section{Test Case Scenarios}

Test case scenarios are a method to fulfill the satisfaction of the user. In this process of testing we have to test each and all the modules or components of the system to get the result, whether the result might be correct or not correct. If the result is not correct or the output is not correct as the user wants, then the process of testing repeats and now the system is tested very deeply through all parts, all functionalities. Each module of the system is tested separately in order to check whether it works according to user requirements or not.

\subsection{Fee Comparison}

When users visit the website they check the fee schedule of all the universities because it is the most important factor for all those who want to get admission in university. Basically in this web portal system all the universities' fees are mentioned, and the main purpose is to give an easy way to the user with just one click. The web portal removes the difficulty of the user having to visit many different university websites to check the fee of a specific course. Users simply type any course name and then they will see all the fee comparison about that selected course in all universities offering that course.

\textbf{Execution Procedure:}
In this module we have checked the execution of fee structure. The main screen is accessed by the URL of the website; then different services of the system are shown in the main screen. When the user wants to see any fee structure they click on the fee structure button; then a new page is opened which includes four search fields (select university, select category, select department, select course) and one search button. After selection, the system will extract data from different sites and after filtration store it in the database in a defined format. After being stored in the database, the output is shown on the screen in the form of a div. If the user wants to compare with others, a checkbox is clicked and then all comparisons are shown.

\textbf{Description:} Fee comparison makes clear the fee comparison of different universities. It shows different universities' fee structures in the form of a comparison.

\textbf{Objectives:}
\begin{itemize}
    \item It helps the user to know about fees, whether it is expensive or not.
    \item It is easy to find the fee structure of universities.
\end{itemize}

\textbf{Test Driver:}
\begin{itemize}
    \item Dreamweaver, XAMPP.
    \item This system needs only a computer; no need of any other specific device.
    \item This application can run on PCs, laptops and mobile devices.
\end{itemize}

\subsection{Program Offered Module}

It is another major component of the system because when a user wants to get information about different universities, they first check the programs that which university offers, because every university does not offer all the programs. Some universities are only engineering universities so they just offer engineering related programs; some are medical universities so they just offer medical related programs. After checking this module, users then select the programs which are required.

\textbf{Execution Procedure:}
In this module we have checked the execution of program offered. The main screen is accessed by the URL of the site; then when the user wants to see any program offered in different universities they click on the program offered button, then a page is opened which includes four search fields (select category, select department, select course, location) and one search button. After selection, the system will extract data from different websites and after filtration store in the database in a defined format. After being stored in the database, the output is shown on the screen in the form of a table.

\textbf{Description:} Users get information about which programs are offered by universities, helping them to inform about which programs are offered by which university.

\textbf{Objectives:}
\begin{itemize}
    \item Very helpful for users because it shows all those programs which are available in different universities.
    \item Users can choose the right program with better information.
    \item Also, other related programs are linked by the selective programs.
\end{itemize}

\textbf{Test Driver:}
\begin{itemize}
    \item Dreamweaver, XAMPP.
    \item This application can run on PCs, laptops and mobile devices.
\end{itemize}

\subsection{Searching Module}

The search module is very helpful for users because by using it, they can easily access or search for the information that they want. Users can search a university by its name, search by any course name and also search by location.

\textbf{Execution Procedure:}
In this module we have checked the execution of the search mechanism. The main screen is accessed by the URL of the website; then when the user wants to search any information in different universities they go to the search button, then a page is opened which includes four search fields (select category, select department, select course, location) and one search button. After selecting, the system will extract data from different sites and after filtration, store data in the database in a defined format.

\textbf{Description:} Users can search universities, programs and cities.

\textbf{Objectives:}
\begin{itemize}
    \item It gives proper information to users about the system and how to select a university.
    \item Not difficult; very easy to understand.
\end{itemize}

\textbf{Test Driver:}
\begin{itemize}
    \item Dreamweaver, XAMPP.
    \item This application can run on PCs, laptops and mobile devices.
\end{itemize}

\subsection{Ranking}

Basically, ranking helps the user to understand the performance of universities. Many users search universities by the ranking given by HEC because some students with good grades want to select a highly ranked university. If a user wants to get admission in a high-ranked university, they can get the complete information about ranking.

\textbf{Execution Procedure:}
In this module we have checked the execution of Ranking. The main screen is accessed by the URL of the site; then when the user wants to see the ranking of different universities, they click on the ranking button. Then a page is opened which includes two search fields (select university, select department) and one search button; after being selected, the system will extract data from different sites and after filtration store in the database in a defined format.

\textbf{Description:} Users can easily know about the rank of universities which are given by HEC.

\textbf{Objectives:}
\begin{itemize}
    \item A proper way to know about universities' rank.
    \item It is good to know about the progress of universities.
\end{itemize}

\textbf{Test Driver:}
\begin{itemize}
    \item Dreamweaver, XAMPP.
    \item This application can run on PCs, laptops and mobile devices (internet connection is required).
\end{itemize}

\subsection{Faculty}

It is necessary to know about the faculty of the university in which the user wants to get admission. Faculty helps the user about their program relating because in faculty there is all information about teachers such as their qualification, teaching experience, behavior in class, designation and other profile information.

\textbf{Execution Procedure:}
In this module we have checked the execution of faculty. The main screen is accessed by the URL of the site; then when the user wants to see faculty of different universities, they click on the faculty button. Then a page is opened which includes three search fields (select university, select department, location) and one search button; after being selected, the system extracts data from different sites and after filtration stores it in the database in a defined format.

\textbf{Description:} Users can easily check the faculty.

\textbf{Objectives:}
\begin{itemize}
    \item A proper way to know about universities' faculty.
    \item The purpose of faculty is to show how many teachers are in different courses.
    \item How many teachers are highly educated and how many are newly appointed.
\end{itemize}

\textbf{Test Driver:}
\begin{itemize}
    \item Dreamweaver, XAMPP.
    \item This application can run on PCs, laptops and mobile devices (internet connection is required).
\end{itemize}

\subsection{Eligibility Criteria}

When users search different universities, they check the eligibility criteria. Eligibility criteria is an important module for getting admission, and the benefit of this module is that it saves time because users can check the eligibility criteria to understand about getting admission before exploring everything else. All the criteria about admission is mentioned in eligibility criteria.

\textbf{Execution Procedure:}
In this module we have checked the execution of criteria. The main screen is accessed by the URL of the site; then when the user wants to see criteria of any university or different universities, they click on the criteria button. Then a page is opened which includes four search fields (select university, select category, select department, select course) and one search button; after being selected, the system will extract data from different sites and after filtration store in the database in a defined format.

\textbf{Description:} This is the main point of the system because it tells about which students are eligible to get admission and which are not eligible.

\textbf{Objectives:}
\begin{itemize}
    \item It works like a threshold.
    \item It gives proper information to users about getting admission in different courses.
    \item It also shows the university's requirements including any tests conducted by the university or other institutes.
\end{itemize}

\textbf{Test Driver:}
\begin{itemize}
    \item Dreamweaver, XAMPP.
    \item This application can run on PCs, laptops and mobile devices.
\end{itemize}

\subsection{Module Analysis}

Table~\ref{tab:module_analysis} shows the evaluation of different modules related to the proposed system.

\begin{table}[H]
\centering
\caption{Module Analysis}
\label{tab:module_analysis}
\begin{tabularx}{\textwidth}{|c|l|X|X|X|c|}
\hline
\textbf{Sr.} & \textbf{Module} & \textbf{Test Scenario} & \textbf{Expected Result} & \textbf{Actual Result} & \textbf{Status} \\
\hline
1 & Web Crawler & Extract data through unique id & Extract data of websites against the id & Same as expected & Pass \\
\hline
2 & Filter Data & Filter the data & Give the desired data & Same as expected & Pass \\
\hline
3 & Data Store & Filtered data stored in database in desired format & Store extracted data in database & Same as expected & Pass \\
\hline
4 & Retrieve Data & Retrieve data from database & Data retrieved from database on user request & Same as expected & Pass \\
\hline
\end{tabularx}
\end{table}

\section{Testing Result and Analysis}

The basic goal of this web application is to extract data from different university websites and accumulate it at one platform. It gives the easiest way to continue the study by getting admission in a university after knowing all information where the user wants to get admission. This is a very simple web-based system to understand for retrieving information related to academia. The system is fully tested and finally gives the proper result within a few seconds of the user query.
